\chapter{Conclusion}
\label{chap:conclusion}

% =====================================================================
% OUTLINE CHAPTER. Bullet notes mark what each section argues. Write as prose
% once Chapter 5 numbers exist, so contributions and limitations cite results.
% =====================================================================

% WRITE-UP NOTES: open by answering the research question from Section 1.2
% directly, in one or two sentences, grounded in the Chapter 5 results. Avoid a
% hollow recap; say what was shown, then move to contributions.

\section{Contributions}
\label{sec:contributions}

% The defensible novelty claims (what existing tools do NOT already do):
% - Automated LLM classification of unknown tools against a fixed taxonomy,
%   rather than a human deciding per tool, or a fixed rule set run at everything.
% - LLM-assisted payload synthesis that adapts to arbitrary tool schemas,
%   instead of hand-written test cases per tool shape.
% - Dual rule-based + grounded LLM-judge evaluation: catches known-format and
%   semantic leaks, and overrules an un-evidenced "vulnerable" verdict.
% - A deliberate, documented answer to "does the triage step itself trust
%   untrusted tool descriptions?" (the categorizer floor) — the tool's own
%   LLM step is exposed to the injection it detects, and the design addresses it.
% - One taxonomy and one remediation catalog spanning a dynamic and a static
%   path, so a static finding and a live finding are directly comparable.

\section{Limitations}
\label{sec:limitations}

% From the project's own "known limitations" — state them plainly:
% - Confidence scores are the model's self-report, not a calibrated metric; no
%   precision/recall against a large labeled corpus yet (the victim corpus is
%   small and author-built).
% - Timing-based checks (blind SQLi, DoS) are noisy under load / slow networks.
% - The static payload library is illustrative (DAST-level), not exhaustive.
% - In-memory credential store does not survive a restart or multiple workers.
% - No auth on the API itself (matches the "local application" scope).
% - resource_extraction wildcard probe has no automated success indicator.
% - Cross-session tests (memory poisoning, cross-tenant) are single-session in
%   the current executor; auth test lacks a low-privilege session.
% - No OAuth token auto-refresh between discover and attack; OAuth assumes the
%   backend and browser share a machine.
% - Static attribution is co-location (sink inside a tool's body), not data-flow
%   proof that a parameter reaches the sink; non-Python servers are unparsed.

\section{Future Work}
\label{sec:future}

% Roughly priority-ordered for the research:
% - Build and run a precision/recall benchmark against a larger labeled corpus
%   of vulnerable and known-safe servers — turns "seems to work" into a number.
% - Automate the two-session tests so the auth and cross-session categories are
%   actually exercised (removes the two known misses in Chapter 5).
% - Correlate the static and dynamic paths on one target and report where they
%   agree, where static reached what the payload could not, and vice versa.
% - Static analysis: data-flow instead of co-location; parse non-Python
%   (TypeScript) servers.
% - Normalize findings into their own indexed table; add schema migrations.
% - Add authentication to the API if it ever runs beyond the local machine.
% - Expand payload coverage per category (encoding variants, blind techniques).

\section{Final Remarks}
\label{sec:final}

% WRITE-UP NOTES: one short, substantive paragraph. State what the artifact
% establishes about securing MCP servers and why the taxonomy-shared, two-path
% design is the reusable idea beyond this specific tool. No generic summary.


% Limitations:
% - Scores are category-level base scores: an SQL error message and a full
%   table dump both get 9.4. Temporal and environmental CVSS metrics are not
%   used, so the score ignores exploit maturity and the deployment context.
% Future work:
% - Let the grounded judge adjust the C/I/A metrics per finding from the
%   evidence (e.g. raise C when real data leaked), keeping the rest of the
%   base vector, and recompute the score.
% - Evaluate CVSS v4.0 once tooling support is mature.